% Generated draft: requires mathematical and source review.
\documentclass[11pt]{article}
\usepackage[margin=25mm]{geometry}
\usepackage{fontspec}
\setmainfont{Latin Modern Roman}
\usepackage{amsmath,amssymb,mathtools,unicode-math}
\setmathfont{Latin Modern Math}
\usepackage{enumitem,needspace,xurl,hyperref}
\setlength{\parindent}{0pt}
\setlength{\parskip}{5pt}
\setlength{\emergencystretch}{3em}
\newcommand{\N}{\mathbb{N}}
\newcommand{\Z}{\mathbb{Z}}
\newcommand{\Q}{\mathbb{Q}}
\newcommand{\R}{\mathbb{R}}
\newcommand{\C}{\mathbb{C}}
\newcommand{\F}{\mathbb{F}}
\newcommand{\A}{\mathbb{A}}
\newcommand{\PP}{\mathbb{P}}
\newcommand{\E}{\mathbb{E}}
\newcommand{\eps}{\varepsilon}
\newcommand{\dd}{\,\mathrm{d}}
\begin{document}
\section*{UnsolvedMath ID: 20000959}
\textbf{Source problem code:} catalogue entry
\subsection{Definitions and mathematical statement}
Let \(T=\{t_1,t_2,t_3\}\subset\R^2\) consist of three distinct points, allowing collinearity. A copy means a congruent image under a Euclidean isometry, including reflections, without rescaling. The original conjecture is
\[
 \forall T\ \exists c:\R^2\to\{0,1,2\}\quad
 \forall\text{ Euclidean isometries }g:\quad c\text{ is not constant on }g(T).
\]
No measurability restriction on arbitrary colorings is specified. The supplied partial result studies only the cyclic half-open stripe colorings
\[
 c_w(x,y)=\lfloor y/w\rfloor\pmod3,\qquad w>0.
\]

Write \(D=\operatorname{diam}(T)\), \(u_\theta=(\cos\theta,\sin\theta)\), and
\[
 p_i(\theta)=(t_i-t_1)\cdot u_\theta,\quad
 K=\left\lceil\frac{D+w}{3w}\right\rceil.
\]
For integers \(k_1=0\), \(k_2,k_3\in[-K,K]\), put \(r_i(\theta)=p_i(\theta)-3wk_i\) and \(R_{k_2,k_3}(\theta)=\max_i r_i(\theta)-\min_i r_i(\theta)\). The exact certificate reported by the aid is
\[
 c_w\text{ admits a monochromatic congruent copy of }T
 \quad\Longleftrightarrow\quad
 \exists k_2,k_3,\theta:\ R_{k_2,k_3}(\theta)<w.
\]
The offsets belong to one coherent triple, rather than independent pairwise modular choices. A common vertical translation places all residuals in one half-open interval of width \(w\) exactly when the range is strictly smaller than \(w\); equality is excluded. Varying the projection direction accounts for rotations and reflections, while the common translation is eliminated by the range test.

For each offset pair, the global minimum of \(R\) is found among angles where either \(r_i=r_j\) or \((r_i-r_j)'=0\) for some distinct pair. Each of the three pairs gives at most two angles of each type, so at most twelve critical angles suffice. Therefore at most
\[
 12(2K+1)^2
\]
range evaluations decide whether this particular stripe coloring avoids \(T\), in an exact-real comparison convention. This is a finite certificate for a prescribed \(w\), not a general decision algorithm for arbitrary colorings or unspecified noncomputable real input.

The aid also reports a boundary-safe stripe avoidance theorem for every nondegenerate triangle whose smallest angle is at least \(30^\circ\): some suitable stripe width avoids all congruent copies. Its width choice and full special-case proof are not supplied in the summary and require review; no unprovided width formula is inserted.

Finally, a coloring constant on lines in one fixed direction cannot avoid a collinear three-point set, which can be placed on one such line. A strategy requiring every segment of one prescribed positive length to have differently colored endpoints would be a proper three-coloring of that distance graph, which the aid reports impossible. These are obstructions to particular approaches, not refutations of the original three-point conjecture. The report describes the full conjecture as open and supplies no unrestricted coloring construction.
\subsection{Short English statement}
The full conjecture asks to avoid each prescribed three-point configuration with three colors. The supplied result gives an exact finite test for one cyclic stripe family, including coherent offsets and half-open boundaries, and reports a triangle special case. It does not decide the general coloring question or cover collinear sets with line-constant stripes.
\subsection{Sources}
\begin{itemize}
\item Local registry: permanent UnsolvedMath ID \(20000959\), exact public code \texttt{catalogue entry}; \texttt{aim-combinatorics-notes.json}, \texttt{notes[83]}.
\item AIM workshop \emph{Graph Ramsey theory}, section \emph{Euclidean Ramsey sets}, Conjecture 37: \url{https://www.overleaf.com/read/mnvcscjjysvg}.
\item Supplied research summary, ``A finite certificate for three-stripe triangle avoidance,'' attempt 1. No named external theorem citations, width formula for the angle special case, or complete references for the strategy obstructions are supplied; original remarks and literature fields are empty.
\end{itemize}
\end{document}
